% 380_D4_Spezifitaet_En_ch.tex
% Johann Pascher, 29 Sept. 2026
% Why D4: specificity of the carrier against matched comparison lattices
% (open bridge R95, Doc. 190)

\providecommand{\BM}{\textbf{[B]}}
\providecommand{\KM}{\textbf{[K]}}
\providecommand{\SM}{\textbf{[S]}}
\providecommand{\QM}{\textbf{[Q]}}
\providecommand{\SetzM}{\textbf{[SETZUNG]}}

\chapter*{Why \texorpdfstring{$D_4$}{D4}}
\addcontentsline{toc}{chapter}{Why D4}

\begin{abstract}
FFGFT posits the lattice $D_4$ as the carrier of the compact geometry
$T^4/\mathbb{Z}_3$ (Doc.~314) \SetzM. Doc.~190 lists under R95 as an open
bridge that an adversarial specificity test against matched comparison
carriers is missing. This document examines the carrier. Three criteria
fixed in advance are applied to $\mathbb{Z}^4$, $A_4$, $A_2\oplus A_2$, $D_4$ and to
random lattices: the fixed-point-free $\mathbb{Z}_3$ action with nine fixed
points carried in the corpus (Doc.~330), the kissing number, and the lattice
energy as the Epstein zeta function at equal volume of the fundamental cell.
Only the third criterion is independent of FFGFT. The $\mathbb{Z}_3$ condition
excludes $\mathbb{Z}^4$ and $A_4$; among the remaining carriers $D_4$ has the
lowest lattice energy \KM, and the kissing number 24 is consistent with
Doc.~340. $D_4$ is thereby better motivated as a carrier but remains a posit;
R95 is not closed, and the level of the operator $\hat F$ remains open \SM.
Verification script: 13/13.
\end{abstract}

% ============================================================
\section*{Status markers}
% ============================================================
\begin{center}
\begin{tabular}{@{}lp{10cm}@{}}
\textbf{[B]} & Algebraically proved. \\[3pt]
\textbf{[K]} & Numerically confirmed (verification script). \\[3pt]
\textbf{[S]} & Structurally plausible, open. \\[3pt]
\textbf{[Q]} & Taken from the literature, cited. \\[3pt]
\textbf{[SETZUNG]} & Axiomatic posit. \\
\end{tabular}
\end{center}

% ============================================================
\section{The question}
% ============================================================

A posit is motivated rather than arbitrary if the chosen carrier stands out
among comparable alternatives by quantities that were not tuned to the
results. This is exactly what the GAE logic demands, as applied in the HLV
audit: there an ordinary cubic lattice outperformed the carrier under study in
the spectral measure, and the specificity claim failed (Doc.~330). For $D_4$,
R95 records that $5\nmid|\mathrm{Aut}(D_4)|=1152$ is a lattice-specific fact,
but not a proof of specificity.

The test must therefore meet three conditions. The comparison carriers have
the same dimension and the same volume of the fundamental cell. The criteria
are fixed before the calculation. And a criterion meant to establish
specificity must not be a number that originates in FFGFT. Criteria that
express a corpus condition test only compatibility; they are explicitly
marked as such below.

% ============================================================
\section{Comparison carriers and criteria}
% ============================================================

The comparison covers the cubic lattice $\mathbb{Z}^4$ (the winner in the HLV
audit), the root lattice $A_4$, the sum $A_2\oplus A_2$ and $D_4$, each scaled to
unit volume of the fundamental cell, plus random lattices.

\paragraph{K1: $\mathbb{Z}_3$ compatibility (corpus condition).}
The corpus carries $T^4/\mathbb{Z}_3$ with nine fixed points and a $\mathbb{Z}_3$
action with eigenvalues $\{\omega,\omega^2,\omega,\omega^2\}$; the nine fixed points
follow from $|\det(1-g)|=|(1-\omega)(1-\omega^2)|^2=9$ (Doc.~330, likewise
Docs.~361, 364, 365). What is required is therefore a lattice automorphism $g$
of order~3 without eigenvalue~1. K1 expresses this condition and tests
compatibility, not specificity. The search is complete, because every
automorphism maps shortest vectors to shortest vectors and all four bases
consist of shortest vectors.

Which $\mathbb{Z}_3$ is meant is shown by the complete enumeration of the 1152
automorphisms of $D_4$: 80 have order~3. Of these, 64 have a two-dimensional
fixed space and produce fixed tori instead of isolated fixed points, among
them the Dynkin automorphism of triality. The remaining 16 are
fixed-point-free with exactly nine fixed points on $T^4$ \KM. These are the
16 elements named in Doc.~330; the triality that carries the orbifold is thus
this fixed-point-free class, not the Dynkin automorphism.

\paragraph{K2: Kissing number (consistency).}
The number of shortest lattice vectors. It is a property of the lattice alone,
but FFGFT uses it itself ($\mathrm{Kissing}(D_4)=24$ in Doc.~340). K2 therefore
tests consistency, not specificity.

\paragraph{K3: Lattice energy.}
The Epstein zeta function
\begin{equation}
Z_\Lambda(\nu)=\sum_{v\in\Lambda\setminus\{0\}}|v|^{-\nu}
\end{equation}
at unit volume of the fundamental cell, for $\nu=5,6,8$. It is at the same
time the spectral zeta function of the flat torus of the dual lattice, hence
an operator quantity of the Laplacian, and it contains no free parameter. It
is computed with EpsteinLib (Buchheit et al.) \QM; the library reproduces the
closed-form values for $\mathbb{Z}^4$ and $D_4$ to twelve digits \KM.

% ============================================================
\section{Results}
% ============================================================

\paragraph{K1.}
$\mathbb{Z}^4$ and $A_4$ have no automorphism of order~3 without eigenvalue~1;
on them $T^4/\mathbb{Z}_3$ with nine fixed points is impossible \KM. $A_2\oplus A_2$
and $D_4$ have one, each with $|\det(1-g)|=9$ \KM. The winner of the HLV audit
is thus already excluded by the symmetry condition.

\paragraph{K2.}
The kissing numbers are $D_4$: 24, $A_4$: 20, $A_2\oplus A_2$: 12,
$\mathbb{Z}^4$: 8 \KM. Of the two $\mathbb{Z}_3$-compatible carriers, only $D_4$
has the 24 that enters Doc.~340 as $\mathrm{Kissing}(D_4)$. This is
consistency: the corpus computes with the lattice it posits.

\paragraph{K3.}
\begin{center}
\begin{tabular}{lcccc}
\toprule
$\nu$ & $D_4$ & $A_4$ & $A_2\oplus A_2$ & $\mathbb{Z}^4$ \\
\midrule
5 & \textbf{22.860} & 23.123 & 23.703 & 24.531 \\
6 & \textbf{12.583} & 12.919 & 13.699 & 14.830 \\
8 & \textbf{6.830} & 7.282 & 8.457 & 10.246 \\
\bottomrule
\end{tabular}
\end{center}
$D_4$ has the lowest energy for all three exponents \KM. Of 300 general random
lattices none falls below $D_4$ (best $13.15$ at $\nu=6$), and likewise none of
300 random lattices from the $\mathbb{Z}_3$-compatible family (best $12.72$)
\KM. All 200 small perturbations around $D_4$ increase the energy; $D_4$ is a
local minimum \KM.

It is proved that $D_4$ is a local minimum of the Epstein zeta function among
all four-dimensional lattices (Sarnak and Strömbergsson 2006) \QM. Whether it
is also the global minimum is an open conjecture in four dimensions. The 600
random lattices support it but do not prove it.

% ============================================================
\section{Assessment}
% ============================================================

\paragraph{What is shown.}
Among the four-dimensional carriers that admit the $\mathbb{Z}_3$ action with
nine fixed points carried in the corpus, $D_4$ has the lowest lattice energy
\KM. This is the only one of the three quantities that does not originate in
FFGFT; the kissing number confirms consistency only. $D_4$ is thereby better
motivated as a carrier than before, but it remains a motivated posit
\SetzM. The finding does not replace a proof of specificity.

\paragraph{What remains open.}
The GAE logic also asks about the operator: whether $\hat F$ with
$\xi=\lambda_{\min}(\hat F_{D_4})$ (Doc.~327) yields something distinguishable on
the comparison carriers. This is not tested here \SM. Likewise, the random
search is numerical evidence, not a proof of global minimality within the
$\mathbb{Z}_3$ family; it covers the exponents $\nu=5,6,8$.

\paragraph{Consequences for the bookkeeping.}
R95 is not closed. The required operator test is still outstanding, and new
results receive no register entry.

% ============================================================
\section{Summary}
% ============================================================

\begin{center}
\begin{tabular}{>{\raggedright\arraybackslash}p{3.2cm}cccc}
\toprule
Criterion & $D_4$ & $A_2\oplus A_2$ & $A_4$ & $\mathbb{Z}^4$ \\
\midrule
K1: $\mathbb{Z}_3$ with 9 fixed points (corpus) & yes & yes & no & no \\
K2: kissing number (consistency) & \textbf{24} & 12 & 20 & 8 \\
K3: $Z(6)$ at unit volume & \textbf{12.58} & 13.70 & 12.92 & 14.83 \\
\bottomrule
\end{tabular}
\end{center}

Among the carriers that satisfy the symmetry condition, $D_4$ has the lowest
lattice energy, a quantity that does not originate in FFGFT; the kissing
number is consistent with the corpus. The posit of the carrier is thus better
supported, but it remains a posit.

\paragraph{Verification script.} \texttt{2/python/Dok380\_Skripte/pruef\_380\_d4\_spezifitaet.py}
(13/13): check of the library against closed-form values, K1 to K3,
classification of the order-3 automorphisms of $D_4$, random and perturbation
tests.

\paragraph{Sources.}
Buchheit,~A.~A.; Busse,~J.~K.; Gutendorf,~R.: \textit{Computation and
Properties of the Epstein Zeta Function with Applications to Quantum
Systems.} arXiv:2412.16317 (2024); software EpsteinLib,
\url{https://github.com/epsteinlib/epsteinlib}.
Sarnak,~P.; Strömbergsson,~A.: \textit{Minima of Epstein's zeta function and
heights of flat tori.} Invent.\ Math.\ \textbf{165}, 115--151 (2006).
Conway,~J.~H.; Sloane,~N.~J.~A.: \textit{Sphere Packings, Lattices and
Groups.} 3rd ed., Springer (1999).
